% !TeX root = ../../Book1.tex
% 纲要标题：### 15.4 嵌入的延拓与共轭
% 纲要位置：计划纲要.md，第 548 行。

\section{嵌入的延拓与共轭}
\label{b1:sec:1504}


% 纲要内容（仅作写作计划，尚非教材正文）：
%
% * 中间域嵌入的延拓与自同构；代数闭包上的自同构延拓
% * 多项式根的共轭性
% * 为自同构计数准备工具
%
% ---
%

比较两个根时，我们通常不只关心它们满足同一个方程，还关心替换这个根能否同时替换域中的其他元素。嵌入延拓定理保证代数关系可相容地延续；正规性决定延拓后的像是否仍是原域。本节把这两种作用分清，并整理随后计数所需的形式。

\subsection{中间域嵌入的延拓与自同构}
\label{b1:subsec:1504:01}

\begin{corollary}\label{b1:cor:intermediate-embedding-extension}
设 \(K\subseteq E\subseteq L\) 为域扩张塔，并且 \(L/K\) 代数；设 \(\Omega\) 为包含 \(L\) 的 \(K\) 的代数闭包。每个 \(K\) 嵌入 \(\tau:E\rightarrow\Omega\) 均能延拓为 \(K\) 嵌入 \(L\rightarrow\Omega\)。若 \(L/K\) 正规，将该延拓的目标限制为 \(L\) 后，它为 \(L\) 的 \(K\) 自同构。此外，\(\tau\) 总能延拓为 \(\Omega\) 的 \(K\) 自同构。
\end{corollary}
\begin{proof}
\(L/E\) 代数，对它与嵌入 \(\tau\) 使用\cref{b1:thm:embedding-extension}，便得到 \(\widetilde\tau:L\rightarrow\Omega\)。由于它延拓 \(\tau\)，在 \(K\) 上仍为恒等；当 \(L/K\) 正规时，\cref{b1:thm:normal-embedding-criterion}给出 \(\widetilde\tau(L)=L\)。所以将目标限制为 \(L\) 后，才得到 \(L\) 的自同构。

\(\Omega/K\) 代数，故 \(\Omega/E\) 也代数。再次应用嵌入延拓定理，可将 \(\tau\) 延拓为 \(\Omega\rightarrow\Omega\) 的 \(K\) 嵌入。\cref{b1:prop:algebraic-self-embedding}保证它满射，因而是 \(\Omega\) 的 \(K\) 自同构。
\end{proof}
上述结论也允许 \(E/K\) 无限；在代数闭包中，局部选定的共轭根可以相容地延拓到整个域的自同构。

\subsection{多项式根的共轭性}
\label{b1:subsec:1504:02}

\begin{definition}\label{b1:def:field-conjugate-elements}
设 \(\Omega\) 为域 \(K\) 的代数闭包，\(a,b\in\Omega\) 都在 \(K\) 上代数。如果 \(a,b\) 具有相同的 \(K\) 上最小多项式，则称它们为 \(K\) 上的\term[gonge-yuan]{共轭元}[conjugate elements]，也说它们在 \(K\) 上共轭。
\end{definition}
这等价于 \(a,b\) 是同一个 \(K[x]\) 中首一不可约多项式的根。这里的共轭是域论术语，其定义与群中的 \(gxg^{-1}\) 不同；下面通过域的自同构解释它。

\begin{proposition}\label{b1:prop:conjugates-automorphism}
设 \(\Omega\) 为域 \(K\) 的代数闭包，\(a,b\in\Omega\) 都在 \(K\) 上代数。则 \(a,b\) 在 \(K\) 上共轭，当且仅当存在 \(\Omega\) 的 \(K\) 自同构把 \(a\) 送到 \(b\)。若还有中间域 \(K\subseteq L\subseteq\Omega\) 包含 \(a,b\)，且 \(L/K\) 正规，则也存在 \(L\) 的 \(K\) 自同构实现这一替换。
\end{proposition}
\begin{proof}
同一个不可约多项式的两个根由\cref{b1:lem:simple-embedding-extension}给出 \(K(a)\rightarrow K(b)\) 的 \(K\) 同构，先在单扩张上实现所需替换。再应用\cref{b1:cor:intermediate-embedding-extension}，把它延拓为 \(\Omega\) 的 \(K\) 自同构；若 \(L/K\) 正规，也可将它延拓为 \(L\) 的 \(K\) 自同构。反过来，若自同构送 \(a\) 到 \(b\)，则把 \(m_{a,K}(a)=0\) 映过去，得到 \(m_{a,K}(b)=0\)。由于它首一不可约，必为 \(b\) 的最小多项式。
\end{proof}
共轭性依赖基域。例如 \(\sqrt2\) 与 \(-\sqrt2\) 在 \(\mathbb Q\) 上的最小多项式都是 \(X^2-2\)，但在 \(K=\mathbb Q(\sqrt2)\) 上分别为 \(X-\sqrt2\) 和 \(X+\sqrt2\)，不再相同。扩大基域后，可用的系数增加，原来不能区分的根可能被区分。即使在正特征中出现重根，命题谈的也是不同的元素；重数不会产生新的嵌入。

\subsection{嵌入个数与扩张次数}
\label{b1:subsec:1504:03}

保持 \(K\) 的域同态 \(L\rightarrow\Omega\) 全体记为 \(\operatorname{Hom}_K(L,\Omega)\)。域同态必单射，所以此处它就是嵌入集合。保持 \(K\) 的 \(L\) 的自同构全体记为 \(\operatorname{Aut}_K(L)\)，复合使它成为群。
\symidx{\(\operatorname{Hom}_K(L,\Omega)\)：保持 \(K\) 的域嵌入集合}
\symidx{\(\operatorname{Aut}_K(L)\)：保持 \(K\) 的域自同构群}

\begin{proposition}\label{b1:prop:embedding-degree-bound}
设 \(L/K\) 为有限域扩张，\(\Omega\) 为包含 \(L\) 的 \(K\) 的代数闭包。则
\[
1\leq\#\operatorname{Hom}_K(L,\Omega)\leq[L:K].
\]
若 \(L/K\) 还正规，则这些嵌入恰为 \(\operatorname{Aut}_K(L)\) 的元素。
\end{proposition}
\begin{proof}
\cref{b1:thm:embedding-extension}将 \(\operatorname{id}_K\) 延拓到 \(L\)，所以至少有一个嵌入。为证明上界，写 \(L=K(a_1,\ldots,a_r)\)，令 \(E_0=K\)、\(E_i=K(a_1,\ldots,a_i)\)。给定 \(E_{i-1}\) 的一个嵌入，它到 \(E_i\) 的延拓由 \(a_i\) 的最小多项式经系数变换后的不同根决定。该多项式的次数为 \([E_i:E_{i-1}]\)，故\cref{b1:lem:simple-embedding-extension}给出的延拓数至多为这个数。每个 \(E_i\) 的嵌入都限制为一个 \(E_{i-1}\) 的嵌入，不同限制所对应的延拓彼此不同，因此
\[
\#\operatorname{Hom}_K(E_i,\Omega)
\leq [E_i:E_{i-1}]\,\#\operatorname{Hom}_K(E_{i-1},\Omega).
\]
从 \(E_0\) 的唯一嵌入出发，逐层相乘并用\cref{b1:thm:tower-law}，便得 \(\#\operatorname{Hom}_K(L,\Omega)\leq[L:K]\)。正规情形的断言正是\cref{b1:thm:normal-embedding-criterion}：每个嵌入的像都等于 \(L\)。
\end{proof}
这个上界可能严格。取素数 \(p\)，令 \(t\) 在 \(\mathbb F_p\) 上超越，并置
\[
K=\mathbb F_p(t^p),\qquad L=\mathbb F_p(t).
\]
取包含 \(L\) 的 \(K\) 的代数闭包 \(\Omega\)。由\eqref{b1:eq:rational-power-degree}，\([L:K]=p\)，所以 \(t\) 在 \(K\) 上的最小多项式是 \(X^p-t^p\)。在特征 \(p\) 下，
\[
X^p-t^p=(X-t)^p.
\]
它在 \(\Omega\) 中只有一个不同的根 \(t\)。每个 \(K\) 嵌入 \(L\rightarrow\Omega\) 都须把 \(t\) 送到这个根，因而只能是恒等映射。于是
\[
\#\operatorname{Hom}_K(L,\Omega)=|\operatorname{Aut}_K(L)|=1<p=[L:K].
\]
同时，\(L\) 是 \(X^p-t^p\) 在 \(K\) 上的分裂域，故 \(L/K\) 正规。这个例子表明，正规性保证上述嵌入的像仍为 \(L\)，却不保证嵌入数达到扩张次数；重根的重数不会产生新的嵌入。下一章将证明：对有限扩张，如果没有不可分现象，嵌入数便等于扩张次数；若同时正规，这个数就是自同构群的阶。
